\section{Secrets and Credentials}
\label{sec:secrets}

\subsection{How Aspire Generates and Stores Passwords}

On the first \texttt{dotnet run} invocation of the AppHost, Aspire generates a random password for the PostgreSQL \texttt{postgres} superuser and stores it in the .NET User Secrets store associated with the AppHost project's \texttt{UserSecretsId}. On subsequent invocations, the stored password is read back from User Secrets and passed to the container as the \texttt{POSTGRES\_PASSWORD} environment variable. The container uses this value to initialise the database if it does not already exist, or to authenticate against an existing one.

The User Secrets store is located at:

\begin{lstlisting}
# Windows
%APPDATA%\Microsoft\UserSecrets\<UserSecretsId>\secrets.json

# macOS / Linux
~/.microsoft/usersecrets/<UserSecretsId>/secrets.json
\end{lstlisting}

The \texttt{secrets.json} file is outside the repository directory and is therefore never captured by \texttt{git add} regardless of the \texttt{.gitignore} configuration.

\subsection{Inspecting and Overriding Credentials}

To list all stored secrets for the AppHost:

\begin{lstlisting}
cd src/HobbitGame.AppHost
dotnet user-secrets list
\end{lstlisting}

A typical output for a fresh setup:

\begin{lstlisting}
Parameters:postgres-password = g7K#mP2xQn...
\end{lstlisting}

To override the generated password with a known value (useful for connecting external tools that do not read the Aspire environment injection):

\begin{lstlisting}
dotnet user-secrets set "Parameters:postgres-password" "myDevPassword"
\end{lstlisting}

The key format \texttt{Parameters:\{resource-name\}-password} is the convention Aspire uses for all generated parameter values~\cite{aspiredocs2024}.

\subsection{pgAdmin Credentials}

pgAdmin's default email and password are generated separately and appear as environment variables on the \texttt{pgadmin} resource in the Aspire dashboard. They are accessible at \texttt{PGADMIN\_DEFAULT\_EMAIL} and \texttt{PGADMIN\_DEFAULT\_PASSWORD} in the dashboard's resource-detail view. These credentials can also be overridden via \texttt{WithEnvironment} on the pgAdmin resource builder if a fixed value is preferred during development.

\subsection{Common Credential Pitfalls}

Table~\ref{tab:credential-pitfalls} documents the most common credential-related failures and their resolutions.

\begin{table}[H]
\centering
\caption{Common credential pitfalls when setting up Aspire with PostgreSQL.}
\label{tab:credential-pitfalls}
\small
\renewcommand{\arraystretch}{1.15}
\begin{tabularx}{\textwidth}{@{}>{\raggedright\arraybackslash}p{4cm} >{\raggedright\arraybackslash}X >{\raggedright\arraybackslash}p{3.5cm}@{}}
\toprule
\textbf{Symptom} & \textbf{Cause} & \textbf{Resolution} \\
\midrule
\texttt{password authentication failed for user "postgres"} & Container was created with a different password than the one now in User Secrets & \texttt{docker rm -f postgres} then re-run AppHost; new password stored \\
\addlinespace
\texttt{No connection string named "hobbitdb"} & \texttt{AddDatabase} name does not match \texttt{AddNpgsqlDbContext} name & Ensure both strings are identical \\
\addlinespace
Container starts with empty DB on every run & \texttt{WithDataVolume()} omitted or \texttt{WithLifetime(Persistent)} omitted & Add both calls to the Postgres builder chain \\
\addlinespace
pgAdmin shows no registered server & pgAdmin container was recreated after initial registration & Remove pgAdmin container; Aspire re-registers on fresh start \\
\addlinespace
Credentials appear in \texttt{appsettings.json} & Manual connection string written by developer & Remove from \texttt{appsettings.json}; rely solely on Aspire injection \\
\bottomrule
\end{tabularx}
\end{table}

\subsection{What Must Never Enter Source Control}

Three categories of data must not be committed:

\begin{itemize}
  \item \texttt{secrets.json} files from the User Secrets store (never in the repository directory, so this is enforced by location rather than \texttt{.gitignore}).
  \item Any \texttt{appsettings.\{environment\}.json} file containing a \texttt{ConnectionStrings} or \texttt{Password} key.
  \item \texttt{docker-compose.override.yml} files generated by tooling that may embed passwords.
\end{itemize}

The \texttt{.gitignore} template in this repository excludes \texttt{appsettings.*.json} by default; reviewers should verify this exclusion is present before merging a change that touches the application configuration layer.
